\ssr{ВВЕДЕНИЕ}

Автоматический синтаксический анализ предложений является одной из важных задач компьютерной лингвистики~\cite{Poletaev2023ConstituencyTree}. Данный вид анализа позволяет решать задачи генерации текстов, автоматического реферирования, установления авторства текстов, а также исследовать синтаксические конструкции языков~\cite{Zakirova2023Syntax,Zobnin2009ObjectATE}. 

В отечественной корпусной и прикладной лингвистике преобладающим формализмом в области синтаксического анализа являются системы зависимостей (грамматики зависимостей, деревья зависимостей, деревья подчинения), распространению которого способствовала модель <<Смысл $\leftrightarrow$ Текст>> И.~A.~Мельчука~\cite{inshakova2019razvitie}. В частности, разметку данными системами получил синтаксический подкорпус НКРЯ~\cite{ruscorpora}.

Существуют, однако, альтернативные формализмы, такие как системы составляющих и системы синтаксических групп~\cite{gladkij1985sintaksicheskie}. Появление в недавнем времени работ, посвященных созданию корпусов с разметкой систем составляющих~\cite{Grashchenkov2024RuConst} и систем синтаксических групп~\cite{DemidovEvchenko2024}, а также разработке методов преобразования систем зависимостей в системы составляющих~\cite{Poletaev2023ConstituencyTree}, свидетельствует о необходимости автоматизации построения данных видов синтаксических систем, как на основе исходных текстовых данных, так и уже существующих систем зависимостей.

Цель работы~---~разработка метода построения синтаксических деревьев для ограниченного естественного русского языка на основе наиболее употребимых грамматик.

Для достижения поставленной цели необходимо выполнить следующие задачи:
\begin{itemize}
	\item проанализировать предметную область автоматизированного построения синтаксических деревьев ограниченного естественного русского языка;
	\item провести сравнительный анализ существующих методов построения синтаксических деревьев и сформулировать требования к разрабатываемому методу;
	\item формализовать постановку задачи в виде IDEF0-диаграммы;
	\item сформулировать ограничения предметной области;
	\item разработать метод построения синтаксических деревьев ограниченного естественного русского языка;
	\item обосновать выбор программных средств реализации и разработать программное обеспечение реализующее описанный метод;
	\item определить зависимость времени построения синтаксических деревьев от количества словоформ в предложении;
	\item оценить точность и полноту разработанного метода.
\end{itemize}